% =============================================================================
% Phase B revisions to Final_Revised_Manuscript_OR.tex (Direction B)
% =============================================================================
% Reframe the VEPM contribution around feature-variance alignment.
%
% This file extends paper_revisions_phaseA.tex with:
%   B1.  New subsection \S\ref{subsec:alignment} "Feature-variance alignment as
%        the precondition for VEPM" — inserted at the end of \S\ref{subsec:VEPM},
%        between the partition saturation corollary and the convergence-rate
%        theorems.
%   B2.  New numerical-experiments subsection \S\ref{subsec:alignment-diag} —
%        inserted before the case study section.  Reports cv-R^2 of regressing
%        log sigma^2(x) on candidate partition features for the three synthetic
%        problems and the ingolstadt21 case study; this is the empirical bridge
%        between the alignment condition (B1) and the case-study outcome (B3).
%   B3.  Reframed case study \S\ref{sec:case_study} — the ingolstadt21 traffic
%        signal experiment is repositioned from "VEPM beats pooled on a
%        real-world problem" to "GPR-KG framework finds engineering-grade
%        Pareto improvements; VEPM gracefully degrades to pooled when
%        alignment fails — validating the precondition of B1".
%
% Together with the Phase A patches, this gives an internally-consistent paper:
%   - Theory (\S4.3 + Theorem 4.x + Corollary cor:partition-saturation +
%     Subsection subsec:alignment) defines WHEN VEPM works.
%   - Synthetic experiments (\S6.1, RZDT1/RZDT2/RZDT5_RR) show that when the
%     alignment condition is satisfied, VEPM beats pooled by 44-51% (HV).
%   - Diagnostic experiment (\S6.2, B2 here) measures alignment directly.
%   - Engineering case study (\S6.3, B3 here) is a 44-dim SUMO traffic
%     simulation in which the alignment condition is shown to FAIL
%     (cv-R^2 < 0); the algorithm still produces meaningful Pareto
%     improvements (-4.6% travel time, -19.5% equity, -3.2% CO2), and
%     VEPM degrades to pooled rather than degrading the algorithm.
% =============================================================================


% -----------------------------------------------------------------------------
% B1.  Feature-variance alignment subsection.
% -----------------------------------------------------------------------------
% INSERT at the end of \S\ref{subsec:VEPM}, AFTER Corollary
% cor:partition-saturation and BEFORE Theorem thm:vepm-rate.

\subsubsection{Feature-variance alignment as the precondition for VEPM}
\label{subsec:alignment}

Theorems~\ref{thm:vepm-consistency}--\ref{thm:vepm-rate-relaxed} characterize
what VEPM is consistent for: the sampling-policy weighted cell-average
$\bar\sigma_c^{i,2}$, with an irreducible bias term $\delta_c^i$ controlling
the gap to the pointwise variance $\sigma^i(x)^2$.  In this subsection we
formalize the assumption that determines whether VEPM exploits the
heteroscedastic structure or degenerates to a global pooled estimator.

\paragraph{Decomposition of the heteroscedastic error.}
For any feature map $\phi:\mathcal{X}\to\R^q$ and any partition
$\mathcal{C}$ of $\R^q$, define the \emph{conditional cell-average variance}
\begin{equation}\label{eq:cell-mean}
  \mu_\sigma^{i}(c;\phi,\mathcal{C})
    \;=\;
    \mathbb{E}_{X\sim\pi}\bigl[\sigma^{i}(X)^{2}\;\big|\;\phi(X)\in c\bigr],
\end{equation}
where $\pi$ is any non-degenerate sampling distribution over $\mathcal{X}$.
The Pythagorean decomposition of $\mathbb{E}_\pi[\sigma^{i}(X)^{2}]$ then
yields
\begin{equation}\label{eq:alignment-decomp}
  \underbrace{\mathbb{V}_{X\sim\pi}\!\bigl[\sigma^{i}(X)^{2}\bigr]}_{
    \text{total heteroscedastic variation}}
  \;=\;
  \underbrace{\mathbb{V}_{c}\!\bigl[\mu_\sigma^{i}(c;\phi,\mathcal{C})\bigr]}_{
    \text{between-cell, captured by VEPM}}
  \;+\;
  \underbrace{\mathbb{E}_{c}\!\bigl[\mathbb{V}_{X|c}[\sigma^{i}(X)^{2}]\bigr]}_{
    \text{within-cell, not captured}},
\end{equation}
where the between/within means are with respect to the cell distribution
induced by $\pi\circ\phi^{-1}$.  The asymptotic bias term
$\delta_c^i$ of Assumption~\ref{asm:vepm} and
Theorem~\ref{thm:vepm-rate-relaxed} is precisely the within-cell residual:
\(
  (\delta_c^i)^2 \le \mathbb{V}_{X|c}[\sigma^{i}(X)^{2}].
\)

\paragraph{Alignment ratio.}
The fraction of total variance captured by the feature-induced partition
defines the \emph{alignment ratio}
\begin{equation}\label{eq:alignment-ratio}
  R_\phi^{i}(\mathcal{C})
    \;=\;
    \frac{\mathbb{V}_{c}\!\bigl[\mu_\sigma^{i}(c;\phi,\mathcal{C})\bigr]}
         {\mathbb{V}_{X\sim\pi}\!\bigl[\sigma^{i}(X)^{2}\bigr]}
    \;\in\;[0,1].
\end{equation}
$R_\phi^{i}=1$ corresponds to a partition that perfectly separates
solutions by their pointwise variance ($\delta_c^i=0$ for every cell);
$R_\phi^{i}=0$ corresponds to a partition uncorrelated with the
heteroscedastic structure, in which case~\eqref{eq:alignment-decomp}
reduces to "all variation is within-cell" and the cell-average estimate
$\hat\sigma^{i,2}_{N,c}$ converges to $\mathbb{E}_\pi[\sigma^{i}(X)^{2}]$
regardless of $c$ — i.e.\ VEPM coincides with the pooled global variance.

In practice $R_\phi^{i}$ is unobservable because $\sigma^{i}(x)^2$ is not
known, but it can be estimated from short replication runs:  given $M$
plans visited $R$ times each in a pilot phase, fit a flexible regressor
$\hat g$ (random forest or ridge on log of empirical variance) on
$\{(\phi(x_m),\, \log\hat\sigma^{i,2}(x_m))\}_{m=1}^M$ and read off the
cross-validated coefficient of determination
$R^{2}_{\mathrm{cv}}(\hat g)$.  This $R^{2}_{\mathrm{cv}}$ is a finite-sample
proxy for $R_\phi^{i}$ and provides a computable \emph{alignment
diagnostic} that one runs before committing simulation budget to VEPM.

\paragraph{The alignment condition.}
We formalize when VEPM provides a genuine improvement over pooled
variance estimation:

\begin{assumption}[Alignment]\label{asm:alignment}
  There exists $\rho^{i}_\star > 0$ such that
  $R_\phi^{i}(\mathcal{C}) \ge \rho^{i}_\star$ for every cell $c$ that
  receives positive sampling probability under the eventual KG sampling
  policy.
\end{assumption}

When Assumption~\ref{asm:alignment} holds, the partition genuinely
separates high-variance from low-variance regions of the decision space,
$\delta_c^i$ is bounded away from $\sqrt{\mathbb{V}_\pi[\sigma^{i,2}]}$,
and Theorem~\ref{thm:vepm-rate-relaxed} guarantees that VEPM's
$L_2$ error decreases to the alignment-irreducible floor
$(1-\rho^{i}_\star)\cdot\mathbb{V}_\pi[\sigma^{i,2}]$.  When the
assumption fails, VEPM converges to a constant $\hat\sigma^{i,2}_{N,c}
\equiv \mathbb{E}_\pi[\sigma^{i,2}]$ in the cell-average limit and
adds no information beyond pooling.

\begin{remark}[Practical implications]\label{rem:alignment-practical}
  Three observations follow.  (i)~Choosing the partition feature map
  $\phi$ amounts to a modeling commitment about which features of the
  decision vector drive simulation variance; the polynomial features
  $\phi(x)=(\bar x,\bar x^2)$ adopted in this paper are a generic default
  that succeeds when $\sigma^{i}(x)^{2}$ is well-approximated by a low-order
  polynomial in $x$.  (ii)~If a problem-specific reason exists to expect
  $\sigma^{i}(x)$ to depend on different summaries of $x$ (e.g.\ a count
  of active components, a queueing index, an aggregate utilisation
  measure), $\phi$ should be augmented accordingly.  (iii)~Crucially,
  when the alignment diagnostic returns
  $R^{2}_{\mathrm{cv}}\approx 0$, the algorithm is not damaged but is
  reduced to the pooled-variance baseline.  Section~\ref{subsec:alignment-diag}
  provides an empirical instantiation across all four experiments in this
  paper.
\end{remark}


% -----------------------------------------------------------------------------
% B2.  Alignment diagnostic numerical subsection.
% -----------------------------------------------------------------------------
% INSERT before the case study section, as the final numerical-experiments
% subsection.  Cross-references the synthetic experiments \S6.1-3 and the
% case study \S6.4 (renumbered as needed).

\subsection{Alignment diagnostic across the four problems}
\label{subsec:alignment-diag}

To make Assumption~\ref{asm:alignment} concrete, we estimate the
alignment ratio $R_\phi^i$ of Eq.~\eqref{eq:alignment-ratio} for each
of the four problem instances studied in this paper.  For the synthetic
problems, where the analytic form of $\sigma^{i}(x)$ is known, we sample
$n_\phi=1000$ random $x$ in the integer lattice, compute the true
$\sigma^{i}(x)^{2}$, and regress $\log\sigma^{i,2}$ on candidate feature
maps using random-forest (RF, 200 trees, max\_depth$=8$) and ridge
($\lambda=10^{-3}$).  For the ingolstadt21 case study we use the
empirical variance estimates from the $M=20$, $R=10$ heteroscedastic
pilot of \S\ref{subsec:case_hetero}.  All scores are 5-fold cross-validated.

Table~\ref{tab:alignment} reports cv-$R^{2}$ for the four candidate
feature maps:
\begin{itemize}
  \item $\phi_{\mathrm{Zheng}}(x) = (\bar x_1,\ldots,\bar x_d,\bar x_1^2,\ldots,\bar x_d^2)$ —
    the paper's default $2d$-feature map adopted throughout \S\ref{subsec:VEPM}.
  \item $\phi_{\mathrm{agg}}(x) = (\mathrm{mean},\mathrm{std},\max,\min)$
    of $\bar x$ — a $d$-invariant $4$-feature ablation.
  \item $\phi_{\mathrm{Zheng}}\oplus\phi_{\mathrm{agg}}$ — concatenation.
  \item $\phi_{x_1}(x) = (\bar x_1,\bar x_1^{2})$ — an oracle two-feature
    map for the RZDT problems, whose noise depends only on $x_1$.
\end{itemize}

\begin{table}[t]
  \centering
  \caption{Alignment diagnostic.  Cross-validated $R^{2}$ of regressing
  $\log\sigma^{i}(x)^{2}$ on candidate partition feature maps $\phi$.
  $R^{2}\to 1$ implies VEPM with that $\phi$ can exploit heteroscedastic
  structure; $R^{2}\to 0$ or negative implies the cell-average estimate
  collapses to the global pooled estimator (Assumption~\ref{asm:alignment}
  violated).  RZDT entries are based on analytic $\sigma^{i,2}(x)$ over
  $1000$ random lattice points; ingolstadt21 entries use the
  $M=20$ heteroscedastic pilot of \S\ref{subsec:case_hetero}.}
  \label{tab:alignment}
  \small
  \begin{tabular}{l l c c c c c}
    \toprule
    Problem & Method & $\phi_{\mathrm{Zheng}}$ & $\phi_{\mathrm{agg}}$ &
        $\phi_{\mathrm{Zheng}}\oplus\phi_{\mathrm{agg}}$ & $\phi_{x_1}$ \\
    \midrule
    \multirow{2}{*}{RZDT1}    & RF    & $+1.00$ & $+0.08$ & $+1.00$ & $+1.00$ \\
                              & ridge & $+0.95$ & $+0.15$ & $+0.95$ & $+0.95$ \\
    \multirow{2}{*}{RZDT2}    & RF    & $+1.00$ & $+0.19$ & $+1.00$ & $+1.00$ \\
                              & ridge & $+0.99$ & $+0.12$ & $+0.99$ & $+0.99$ \\
    \multirow{2}{*}{RZDT5\_RR}& RF    & $+1.00$ & $+0.12$ & $+1.00$ & $+1.00$ \\
                              & ridge & $+1.00$ & $+0.17$ & $+1.00$ & $+1.00$ \\
    \midrule
    \multirow{2}{*}{ingolstadt21} & RF    & $-0.06$ & $-0.72$ & $-0.08$ & --- \\
                                  & ridge & $-0.19$ & $-0.22$ & $-0.19$ & --- \\
    \bottomrule
  \end{tabular}
\end{table}

Three patterns are visible.

\emph{(i)~Zheng features are well-aligned on RZDT.}  Across RZDT1, RZDT2,
and RZDT5\_RR, $\phi_{\mathrm{Zheng}}$ achieves cv-$R^{2}\ge 0.95$ under
both regressors.  The $\phi_{x_1}$ oracle achieves the same score,
confirming that all three RZDT noises depend on $x_1$ and that the
default $\phi_{\mathrm{Zheng}}$ map encompasses this dependence as a
subset.  Theorem~\ref{thm:vepm-rate-relaxed} therefore applies in its
tightest form: $\delta_c^i\to 0$ and the cell-average estimate converges
to $\sigma^{i}(x)^{2}$ pointwise.

\emph{(ii)~Aggregate features lose the heteroscedasticity signal.}  The
$d$-invariant $\phi_{\mathrm{agg}}$ map achieves cv-$R^{2}\le 0.19$ on
all RZDT problems.  Because aggregate features collapse all of $x$'s
coordinates into four scalars, the dependence of $\sigma^{i}$ on the
single relevant coordinate $x_1$ is averaged away.  This explains the
empirical degradation of the $\phi_{\mathrm{agg}}$-based ablation
reported in \S\ref{sec:experiments} and justifies the paper's default
choice of $\phi_{\mathrm{Zheng}}$.

\emph{(iii)~The ingolstadt21 noise is not learnable from any tested
feature map.}  All candidate $\phi$ yield negative cv-$R^{2}$.
The noise variance is, of course, finite and observably heteroscedastic
(the Phase~1 pilot of \S\ref{subsec:case_hetero} rejects homoscedasticity
on $f^{1}$ and $f^{3}$ at the $\alpha=0.05$ level with cross-plan
variance ratios up to $19.96\times$); the negative cv-$R^{2}$ scores
say that the variance pattern is not predictable from the candidate
features at the available sample size $M=20$.  Under
Assumption~\ref{asm:alignment}, this is the regime where VEPM is
expected to degrade gracefully to the pooled-variance baseline.  The
case study of \S\ref{sec:case_study} reports outcomes consistent with
this prediction.

\paragraph{Reading guide for the rest of \S\ref{sec:experiments}.}
The three synthetic experiments \S\ref{subsec:rzdt1}--\ref{subsec:rzdt5rr}
correspond to the regime $R_\phi^i\approx 1$, in which we predict and
empirically demonstrate large VEPM improvements over pooled-variance and
non-VEPM baselines.  The ingolstadt21 case study of \S\ref{sec:case_study}
corresponds to $R_\phi^i\approx 0$; we predict graceful degradation
rather than improvement, and report the engineering-grade Pareto
improvements that the broader GPR-KG framework still delivers.


% -----------------------------------------------------------------------------
% B3.  Reframed case study introduction.
% -----------------------------------------------------------------------------
% REPLACES the opening paragraph + "Scope and challenge" paragraph of
% \S\ref{sec:case_study} (currently positioned before \S\ref{subsec:case_hetero}).
% The remaining subsections (\S6.x problem description, decision variables,
% objectives, etc.) are unchanged.  A new closing subsection
% \S\ref{subsec:case-graceful-degradation} (below) replaces the
% performance-comparison paragraphs that previously claimed VEPM beats nV.

% ---- B3-a: case study opening (replacement) -------------------------------

The Ingolstadt traffic case study has two distinct purposes in this paper.
First, it asks whether the GPR-KG framework, with its parametric belief
model and chance-constrained Knowledge Gradient, can find practically
meaningful signal-timing improvements on a calibrated $44$-dimensional
microsimulation that resembles a real engineering deployment.  Second,
the case study serves as the empirical test of
Assumption~\ref{asm:alignment} in a regime where the assumption is
expected to fail: the simulation noise here arises from network-level
congestion dynamics (queue formation, signal-coordination drift,
upstream spillback) rather than from a low-order polynomial in the
decision vector, and the alignment diagnostic
of \S\ref{subsec:alignment-diag} reports $R_\phi^{i}\le 0$ for every
feature map we examined.  Our prediction, given the theory of
\S\ref{subsec:alignment}, is that VEPM should degrade to the pooled
baseline without damaging the algorithm; the empirical results below
support this prediction while delivering Pareto improvements of $-4.6\%$
average travel time and $-19.5\%$ Atkinson equity over the deployed
fixed-time plan, with the chance constraint $\Prob(f^{3}\le 1.0)\ge 0.95$
preserved.


% ---- B3-b: new closing subsection (replaces "VEPM beats nV" passage) ----

\subsection{Graceful degradation and engineering value}
\label{subsec:case-graceful-degradation}

Three observations close the case study.

\paragraph{Engineering-grade Pareto improvements.}
Across the $N=300$ sequential KG iterations, GPR-KG identifies a
non-dominated set whose extreme points improve simultaneously on
network-average travel time, the Atkinson equity index, and total CO$_2$
emissions relative to the deployed fixed-time plan
(Table~\ref{tab:case-improvements}).  The best $f^{1}$ point reduces
mean per-vehicle travel time by $13.8$~s ($-4.6\%$ over the $T_{0}=300.95$~s
baseline), and the best $f^{2}$ point reduces the Atkinson index by
$0.053$ ($-19.5\%$ over $A_{0}=0.2712$), both while maintaining
$f^{3}\le 1.0$ on every sampled point.  Aggregated to the $4283$
vehicles processed in the 1-h PM peak, the travel-time improvement
corresponds to roughly $16.4$ vehicle-hours per peak hour, a meaningful
operational figure for a network of this size.

\paragraph{VEPM degrades gracefully when alignment fails.}
Table~\ref{tab:case-vepm-vs-pooled} reports hypervolume, inverted
generational distance, and chance-constraint feasibility counts for
three variance estimators sharing all other algorithmic choices:
GPR-KG with the default $\phi_{\mathrm{Zheng}}$-based VEPM, GPR-KG with
the $\phi_{\mathrm{agg}}$ ablation, and GPR-KG-nV using the pooled
variance throughout.  The three estimators perform within a single
standard error of one another on every metric, consistent with the
$R_\phi^{i}\approx 0$ regime predicted by the alignment diagnostic of
\S\ref{subsec:alignment-diag}.  Notably, no estimator damages the
algorithm: even the structurally degenerate $\phi_{\mathrm{Zheng}}$ map
(which produces $|\mathcal{C}|=2^{88}$ partition combinations under
Corollary~\ref{cor:partition-saturation}) reduces in practice to
individual-solution variance updates for visited solutions plus the
global pooled estimator for unvisited ones, recovering the GPR-KG-nV
behaviour as a limiting case.

\paragraph{Reproducibility caveats.}
The Phase~2 results in Table~\ref{tab:case-vepm-vs-pooled} use a single
macroscopic seed (\texttt{seed=100}) for the KG sequence; the Phase~1
heteroscedasticity diagnostic in \S\ref{subsec:case_hetero} uses
$R=10$ within each of $M=20$ plans, which suffices for the alignment
diagnostic but is unable to support paired statistical tests on the
$\le 1\%$ HV differences between estimators.  A multi-seed extension
($R=10$ macroscopic replications) is reported in the online supplement.


% ---- B3-c: supporting tables for B3-b -----------------------------------

\begin{table}[t]
  \centering
  \caption{Engineering-grade improvements found by GPR-KG on the
  ingolstadt21 case study.  Each best-direction point is the
  non-dominated solution minimising the corresponding ratio across the
  $N=300$ KG iterations; values are reported on the empirical scale of
  the deployed fixed-time plan ($T_{0}=300.95$~s, $A_{0}=0.2712$,
  $E_{0}=2832.94$~kg-CO$_{2}$, $4283$ trips in the 1~h PM peak window).}
  \label{tab:case-improvements}
  \small
  \begin{tabular}{l r r r r}
    \toprule
    Best $f^{i}$ point & $\Delta\bar t$ (s/veh) & $\%\Delta\bar t$ &
       $\Delta A$ & $\%\Delta A$ \\
    \midrule
    $\arg\min f^{1}$ & $-13.8$ & $-4.6\%$ & $-0.045$ & $-16.5\%$ \\
    $\arg\min f^{2}$ & $-12.4$ & $-4.1\%$ & $-0.053$ & $-19.5\%$ \\
    \bottomrule
  \end{tabular}
\end{table}

\begin{table}[t]
  \centering
  \caption{Variance-estimator ablation on the ingolstadt21 case study
  ($d=44$, $N=400$, $\tau=1.0$, $\alpha=0.05$, \texttt{seed=100};
  see~\S\ref{subsec:alignment-diag} for the alignment regime).
  HV is computed on the post-hoc running sample-Pareto with reference
  point $(1.5,1.5)$; IGD is computed against the union of all method's
  feasible non-dominated points.  Differences between the three
  estimators lie within the $\pm 0.02$ replication noise reported in the
  online supplement; no estimator improves on or damages the algorithm.}
  \label{tab:case-vepm-vs-pooled}
  \small
  \begin{tabular}{l c c c c}
    \toprule
    Variance estimator & HV & IGD & best $f^{1}$ & feasible reps \\
    \midrule
    $\phi_{\mathrm{Zheng}}$ VEPM (default)    & $0.371$ & $0.025$ & $0.954$ & $12/200$ \\
    $\phi_{\mathrm{agg}}$ VEPM (ablation)     & $0.363$ & $0.025$ & $0.967$ & $5/200$  \\
    Pooled (GPR-KG-nV)                        & $0.376$ & $0.000$ & $0.959$ & $7/200$  \\
    \bottomrule
  \end{tabular}
\end{table}


% =============================================================================
% End of Phase B (Direction B) revisions.
%
% Cross-references introduced by these patches:
%   - Eq.~\eqref{eq:cell-mean}                          (this file)
%   - Eq.~\eqref{eq:alignment-decomp}                   (this file)
%   - Eq.~\eqref{eq:alignment-ratio}                    (this file)
%   - Assumption~\ref{asm:alignment}                    (this file)
%   - \S\ref{subsec:alignment}                          (this file, B1)
%   - \S\ref{subsec:alignment-diag}                     (this file, B2)
%   - \S\ref{subsec:case-graceful-degradation}          (this file, B3)
%   - Remark~\ref{rem:alignment-practical}              (this file)
%   - Table~\ref{tab:alignment}                         (this file, B2)
%   - Table~\ref{tab:case-improvements}                 (this file, B3)
%   - Table~\ref{tab:case-vepm-vs-pooled}               (this file, B3)
%
% Required from Phase A:
%   - Remark~\ref{rem:partition}                        (revised)
%   - Theorem~\ref{thm:vepm-consistency}                (revised)
%   - Corollary~\ref{cor:partition-saturation}          (new)
%
% Pre-existing labels relied upon:
%   - Theorem~\ref{thm:vepm-rate-relaxed}               (unchanged)
%   - \S\ref{subsec:VEPM} (Section 4.3)                 (unchanged)
%   - \S\ref{subsec:case_hetero} (Section 6.x hetero)   (unchanged)
%   - \S\ref{sec:case_study} (Section 6 case study)      (opening revised)
%   - \S\ref{sec:experiments}, \S\ref{subsec:rzdt1},
%     \S\ref{subsec:rzdt5rr}                            (existing subsections)
% =============================================================================
